\chapter{Stundenaufwand}

\section*{Projektplan (Soll-Zustand)}

Im Projektplan wurden 20 Features mit insgesamt \textbf{160 Stunden} eingeplant (Tabelle~\ref{tab:features}).

\begin{longtable}{>{\bfseries}p{0.5cm} p{13.5cm} r}
  \toprule
  \# & Feature / Aufgabe & Std. \\
  \midrule
  \endfirsthead
  \toprule
  \# & Feature / Aufgabe & Std. \\
  \midrule
  \endhead
  \bottomrule
  \endfoot
  \bottomrule
  \endlastfoot

  1  & Hardware-Auswahl \& Beschaffung (MCU, Sensoren, Solar, Akku, Gehäuse) & 8 \\
  \rowcolor{rowgray}
  2  & Schaltplan \& Aufbau der Wetterstation-Elektronik & 9 \\
  3  & Sensor: Temperatur, Luftfeuchtigkeit \& Luftdruck & 2 \\
  \rowcolor{rowgray}
  4  & Sensor: UV-Index \& Lichtstärke & 2 \\
  5  & Sensor: Partikelkonzentration (SPS30) & 4 \\
  \rowcolor{rowgray}
  6  & Sensor: CO${}_2$-Konzentration (SCD41) & 4 \\
  7  & Sensor: Umgebungslautstärke (MEMS-Mikrofon) & 3 \\
  \rowcolor{rowgray}
  8  & Standalone-Test: Sensoren \& LoRaWAN ohne Systemanbindung & 8 \\
  9  & PCB-Design, Bestellung \& Bestückung & 9 \\
  \rowcolor{rowgray}
  10 & Autarke Stromversorgung: Solar-Laderegler \& LiPo-Akku-Management & 13 \\
  11 & Deep-Sleep \& Power-Management Firmware & 7 \\
  \rowcolor{rowgray}
  12 & LoRaWAN-Firmware: Payload-Codec, Sendezyklus, Duty-Cycle-Management & 12 \\
  13 & Helium-Netzwerk: Registrierung, Console-Integration, Decoder & 5 \\
  \rowcolor{rowgray}
  14 & Backend: LoRaWAN-Datenpipeline zu HydroNode (Amazon SNS) & 8 \\
  15 & Website: Echtzeit-Dashboard für alle Stationssensoren & 14 \\
  \rowcolor{rowgray}
  16 & App: Historische Verlaufsdiagramme (alle Sensorkanäle) & 6 \\
  17 & App: Stations-Verwaltung (Hinzufügen, Benennen, Statusanzeige) & 4 \\
  \rowcolor{rowgray}
  18 & KI-Evaluation: Datensatz-Aufbereitung \& Modellauswahl & 8 \\
  19 & KI-Evaluation: Testing, Benchmarking \& Ergebnisdokumentation & 8 \\
  \rowcolor{rowgray}
  20 & Feldtest, Integrationstests \& Abschlussdokumentation & 16 \\
  \midrule
  \multicolumn{2}{r}{\textbf{Gesamt}} & \textbf{160} \\
\end{longtable}
\vspace{-1em}
\begin{center}
  \small\textit{Tabelle: Geplante Features mit Aufwandsschätzung aus dem Projektplan (Pitch, 26.03.2026)}
  \label{tab:features}
\end{center}

\section*{Soll/Ist-Vergleich}

Tabelle~\ref{tab:vergleich} zeigt den Vergleich zwischen dem geplanten und dem tatsächlichen Status für alle Features.

\begin{longtable}{>{\bfseries}p{0.5cm} p{6.5cm} r r l}
  \toprule
  \# & Feature & Geplant & Ist (est.) & Status \\
  \midrule
  \endfirsthead
  \toprule
  \# & Feature & Geplant & Ist (est.) & Status \\
  \midrule
  \endhead
  \bottomrule
  \endfoot
  \bottomrule
  \endlastfoot

  1  & Hardware-Auswahl \& Beschaffung & 8\,h & 7\,h & \status{Fertig} \\
  \rowcolor{rowgray}
  2  & Schaltplan \& Elektronik-Aufbau & 9\,h & 4\,h & \status{Fertig} \\
  3  & Sensor: Temperatur / Feuchte / Druck & 2\,h & 3\,h & \status{Fertig} \\
  \rowcolor{rowgray}
  4  & Sensor: UV-Index (LTR390 statt VEML6075) & 2\,h & 2\,h & \status{Fertig} \\
  5  & Sensor: Feinstaub SPS30 & 4\,h & 3\,h & \status{Fertig} \\
  \rowcolor{rowgray}
  6  & Sensor: CO${}_2$ SCD41 & 4\,h & 3\,h & \status{Fertig} \\
  7  & Sensor: Lautstärke (MAX4466) & 3\,h & 0\,h & \status{Weggelassen} \\
  \rowcolor{rowgray}
  8  & Standalone-Tests Sensoren \& LoRaWAN & 8\,h & 3\,h & \status{Fertig} \\
  9  & PCB-Design, Bestellung \& Bestückung & 9\,h & \textasciitilde56\,h & \status{Fertig} \\
  \rowcolor{rowgray}
  10 & Autarke Stromversorgung Solar \& Akku & 13\,h & \textasciitilde7\,h & \status{In Arbeit} \\
  11 & Deep-Sleep \& Power-Management & 7\,h & \textasciitilde10\,h & \status{Fertig} \\
  \rowcolor{rowgray}
  12 & LoRaWAN-Firmware & 12\,h & \textasciitilde11\,h & \status{Fertig} \\
  13 & Helium-Netzwerk Registrierung \& Decoder & 5\,h & \textasciitilde3\,h & \status{Fertig} \\
  \rowcolor{rowgray}
  14 & Backend: Pipeline zu HydroNode & 8\,h & \textasciitilde3\,h & \status{Fertig} \\
  15 & Website Echtzeit-Dashboard & 14\,h & 10\,h & \status{Fertig} \\
  \rowcolor{rowgray}
  16 & App: Verlaufsdiagramme & 6\,h & \textasciitilde6\,h & \status{Fertig} \\
  17 & App: Stationsverwaltung & 4\,h & \textasciitilde6\,h & \status{Fertig} \\
  \rowcolor{rowgray}
  18 & KI-Eval: Datensatz \& Modellauswahl & 8\,h & \textasciitilde2\,h & \status{In Arbeit} \\
  19 & KI-Eval: Testing \& Benchmarking & 8\,h & 0\,h & \status{Geplant} \\
  \rowcolor{rowgray}
  20 & Feldtest \& Abschlussdokumentation & 16\,h & \textasciitilde7\,h & \status{In Arbeit} \\
  \midrule
  \multicolumn{2}{l}{\textbf{Gesamt (inkl.\ 3D-Gehäuse, nicht geplant)}} & \textbf{160\,h} & \textbf{154\,h est.} & \\
\end{longtable}
\vspace{-1em}
\begin{center}
  \small\textit{Tabelle: Soll/Ist-Vergleich aller Features. Ist-Stunden sind Schätzungen aus dem Arbeitstagebuch (Stand: Zwischenbericht 30.05.2026).}
  \label{tab:vergleich}
\end{center}

\noindent\textbf{Hinweis:} Die Ist-Stunden pro Feature sind Schätzungen. Im Arbeitstagebuch wird nach Datum und Tätigkeit erfasst, nicht nach Feature-Nummer. Deshalb lassen sich die Stunden nur näherungsweise zuordnen.

\section*{Stundenvergleich nach Bereichen}

\begin{table}[H]
  \centering
  \begin{tabularx}{\textwidth}{X r r r}
    \toprule
    \textbf{Bereich} & \textbf{Geplant (h)} & \textbf{Ist (h, est.)} & \textbf{Differenz} \\
    \midrule
    Hardware, Beschaffung, Schaltplan & 17 & \textasciitilde9 & $-8$ \\
    Sensorik (5 von 6 Sensoren) & 12 & \textasciitilde11 & $-1$ \\
    Standalone-Tests & 8 & \textasciitilde3 & $-5$ \\
    PCB-Design, Bestellung, Inbetriebnahme & 9 & \textasciitilde50 & $+41$ \\
    Autarke Stromversorgung (Solar) & 13 & \textasciitilde7 & $-6$ \\
    Firmware (Deep-Sleep, LoRaWAN) & 19 & \textasciitilde21 & $+2$ \\
    Helium-Netzwerk \& Backend & 13 & \textasciitilde6 & $-7$ \\
    Website \& App & 24 & \textasciitilde23 & $-1$ \\
    KI-Evaluation & 16 & \textasciitilde2 & $-14$ \\
    Feldtest \& Dokumentation & 16 & \textasciitilde10 & $-6$ \\
    3D-Gehäuse (nicht geplant) & 0 & \textasciitilde12 & $+12$ \\
    \midrule
    \textbf{Gesamt (exakt laut Arbeitstagebuch)} & \textbf{160} & \textbf{162} & $+2$ \\
    \bottomrule
  \end{tabularx}
  \caption{Stundenvergleich Soll/Ist nach Bereichen (Ist-Werte geschätzt aus Arbeitstagebuch; exakter Gesamtwert laut Arbeitstagebuch)}
  \label{tab:stunden-vergleich}
\end{table}

\paragraph{Wichtigste Beobachtung:} Das PCB-Design hat mit ca.\ 50 Stunden weit mehr Zeit gekostet als die geplanten 9 Stunden. Das entspricht einem Mehraufwand von rund \textbf{+41 Stunden} in diesem Bereich. Dieser Mehraufwand wurde durch Einsparungen in anderen Bereichen teilweise ausgeglichen: Helium und Backend liefen schneller als erwartet, und die Standalone-Tests waren durch frühe Prototypentests bereits abgedeckt. Das 3D-Gehäuse mit ca.\ 12 Stunden war im Projektplan nicht vorgesehen und stellt einen weiteren ungeplanten Mehraufwand dar. Der exakte Gesamtwert von 162 Stunden stammt direkt aus dem Arbeitstagebuch.

\section*{Detailliertes Arbeitstagebuch}

Die folgende Tabelle zeigt den dokumentierten zeitlichen Aufwand auf Basis der erfassten Stunden.

\small
\setlength{\LTleft}{0pt}
\setlength{\LTright}{0pt}
\begin{longtable}{|l|c|l|p{7.5cm}|}
\hline
\textbf{Datum} & \textbf{Stundenanzahl} & \textbf{Kategorie} & \textbf{Beschreibung} \\
\hline
\endfirsthead

\hline
\textbf{Datum} & \textbf{Stundenanzahl} & \textbf{Kategorie} & \textbf{Beschreibung} \\
\hline
\endhead

\hline
\multicolumn{4}{|r|}{\footnotesize Fortsetzung auf der nächsten Seite} \\
\hline
\endfoot

\hline
\endlastfoot

23.03.2026 & 5 & Hardware & Energieeffiziente Kernkomponenten ausgewählt und bestellt (STM32, Laderegler, Konverter, SCD41). \\
\hline
27.03.2026 & 4 & Software & Testprojekt auf Basis der Vorlage aufgebaut, um den Versand von Dummy-Werten stabil zu validieren. \\
\hline
29.03.2026 & 2 & Hardware & Lötarbeiten für die neuen Komponenten (STM32 und BQ25185) durchgeführt und Inbetriebnahme vorbereitet. \\
\hline
30.03.2026 & 4 & Software & Erste Firmware per ST-Link auf den STM32 geflasht und Verbindungsparameter im LoRaWAN-Setup sauber konfiguriert. \\
\hline
30.03.2026 & 2 & Software & Versandzyklen und Low-Power-Ablauf erfolgreich im regulären LoRaWAN-Sendeprozess getestet. \\
\hline
31.03.2026 & 2 & Hardware & Weitere Sensorik für Luftfeuchtigkeit, Temperatur, Druck und UV-Strahlung ausgewählt und nachbestellt. \\
\hline
31.03.2026 & 1 & Hardware & Vergleichende Sendetests mit eigenem und externem Gateway sowie mit unterschiedlichen Antennen durchgeführt. \\
\hline
01.04.2026 & 1 & Hardware & LoRa-Funktionstest innerhalb des Gebäudes K104 durchgeführt. \\
\hline
01.04.2026 & 2 & Software & Projektstruktur überarbeitet, SCD41 integriert und den Build-Prozess über Makefiles verbessert. \\
\hline
03.04.2026 & 2 & Planung & Mindestanforderungen für Batterie- und Solarversorgung der Station rechnerisch ausgelegt. \\
\hline
05.04.2026 & 2 & Software & Energieeffizienz verbessert und den Batteriemonitor LC709203F in die Mess- und Telemetriepfade integriert. \\
\hline
05.04.2026 & 2 & Software & Feldtests bei sehr schwachen Empfangsbedingungen (\textgreater{}10 km) durchgeführt und GitHub-Actions-Tests eingerichtet. \\
\hline
06.04.2026 & 6 & CAD & PCB-Design erstellt sowie Datenblattabgleich und Leistungsbudget-Berechnungen durchgeführt. \\
\hline
07.04.2026 & 6 & CAD & PCB-Layout detailliert ausgearbeitet und für die nächsten Integrationsschritte vorbereitet. \\
\hline
10.04.2026 & 6 & Software & Verbesserungen am Kommunikations- und Applikations-Stack umgesetzt und stabilisiert. \\
\hline
11.04.2026 & 3 & Planung & Vorbereitung auf eigenes PCB-Design und MCU-Wechsel strukturiert abgeschlossen. \\
\hline
13.04.2026 & 6 & Software & SHT45, BMP390, LTR390 und MAX17048 implementiert, eingebunden und funktional verifiziert. \\
\hline
14.04.2026 & 2 & Hardware & Signalstärketests am OTH-Standort durchgeführt; Empfang durch mehrere Stationen erfolgreich nachgewiesen. \\
\hline
15.04.2026 & 6 & Software & Remote-Debugging eingerichtet und zentrale Ablauf- sowie Pre-Wake-Logik gezielt verbessert. \\
\hline
16.04.2026 & 2 & Software & SPS30 in den Messzyklus integriert und in Betrieb genommen. \\
\hline
16.04.2026 & 1 & Software & Manuelle Fan-Cleaning-Abläufe nachgeschärft und zugehörige Fehler behoben. \\
\hline
18.04.2026 & 7 & CAD & PCB- und RF-Design finalisiert und für den Fertigungs- bzw. Integrationsstand freigegeben. \\
\hline
20.04.2026 & 5 & Software & CMake- und Projektstruktur für den STM32WLE5CCU6 neu aufgebaut und migrationssicher eingerichtet. \\
\hline
21.04.2026 & 1 & Planung & LaTeX-Dokumentationsprojekt initial angelegt und Grundstruktur erstellt. \\
\hline
22.04.2026 & 1 & Software & RX-Aktionen in den Kommunikationsablauf integriert. \\
\hline
22.04.2026 & 1 & Software & Fehler im manuellen SPS30-Fan-Cleaning reproduziert und korrigiert. \\
\hline
24.04.2026 & 2 & CAD & PCB-Layout finalisiert und letzte Hardwareanpassungen abgeschlossen. \\
\hline
24.04.2026 & 2 & Software & Firmware gezielt an das finalisierte PCB angepasst und Regressionen geprüft. \\
\hline
26.04.2026 & 6 & Software & Helium SNS in das HydroNodeNetwork integriert und den End-to-End-Datenfluss validiert. \\
\hline
27.04.2026 & 5 & Software & HydroNode-App um zusätzliche Funktionen für HydroNodeStationen erweitert. \\
\hline
27.04.2026 & 3 & Software & Firmware der HydroNodeStation01 an Netzwerkabläufe angepasst und Station in der App eingerichtet. \\
\hline
28.04.2026 & 2 & CAD & PCB finalisiert und Bestellung inkl.\ Bestückung bei JLCPCB aufgegeben. \\
\hline
29.04.2026 & 2 & Software & Energieanalyse des Gesamtsystems ohne Solarladegerät durchgeführt und dokumentiert. \\
\hline
30.04.2026 & 1 & Software & State-of-Charge-Ausgabe des MAX17048 in Messzyklus integriert. \\
\hline
30.04.2026 & 2 & Software & Troubleshooting des SubGHz-Fehlers (write CMD error 1) analysiert und dokumentiert. \\
\hline
01.05.2026 & 3 & Planung & Studiendokumentation erweitert (Hardware, Software, Energieeffizienz, Testing, Open Source). \\
\hline
02.05.2026 & 6 & CAD & 3D-Gehäuse für die Wetterstation in Fusion 360 konstruiert und für den 3D-Druck aufbereitet. \\
\hline
07.05.2026 & 3 & Hardware & Langzeittest ohne Solarversorgung mit dem WioE5-Prototypen gestartet, um den realen Batterieabfall unter Dauerbetrieb zu messen. \\
\hline
11.05.2026 & 4 & Hardware & PCB-Fehlersuche durchgeführt und einen defekten TPS63900-Spannungswandler als Ursache identifiziert und dokumentiert. \\
\hline
12.05.2026 & 3 & Software & Erste Firmware-Uploads auf dem eigenen PCB durchgeführt; Probleme im LoRa-Core analysiert und schrittweise eingegrenzt. \\
\hline
16.05.2026 & 10 & Software & Projektwebsite von Grund auf erstellt und mit Live-Messdaten, Projektbeschreibung und technischen Details befüllt. \\
\hline
18.05.2026 & 2 & CAD & Sensorgehäuse für die Außenmontage modelliert und ersten 3D-Druckdurchlauf gestartet. \\
\hline
20.05.2026 & 3 & CAD & Zweiten Druckdurchlauf des Sensorgehäuses abgeschlossen, Teile nachbearbeitet und auf Passgenauigkeit geprüft. \\
\hline
22.05.2026 & 6 & Software & iOS-App aktualisiert, Fehler behoben und bestehende Webapp-Funktionen erweitert. \\
\hline
26.05.2026 & 4 & Software & Fehler auf den gefertigten PCBs behoben; Firmware erstmals stabil mit externer Batterieversorgung in Betrieb genommen. \\
\hline
27.05.2026 & 1 & Hardware & Kleinere Solarzelle auf Eignung für die Energieversorgung der Station getestet. \\
\hline
28.05.2026 & 3 & Hardware & Outdoor-Gehäuse vollständig montiert, alle Sensoren eingebaut und Feldtest gestartet; Projektwebsite parallel optimiert. \\
\hline
30.05.2026 & 1 & Software & Erste Messwerte aus dem Feldtest ausgewertet und auf Plausibilität sowie Konsistenz geprüft. \\
\hline
06.06.2026 & 2 & Hardware & Plexiglasabdeckung für den Sensorbereich der Außenstation gefertigt und eingepasst. \\
\hline
08.06.2026 & 4 & Software & UV-Index-Berechnung im Firmware-Stack refaktoriert, Sensorkabel überarbeitet und neue Testläufe gestartet. \\
\hline
\textbf{Gesamt} & \textbf{162} &  &  \\
\hline
\caption{Dokumentierter Stundenaufwand nach Datum, Kategorie und Tätigkeit (Stand: 08.06.2026)}\label{tab:stundenaufwand}\\
\end{longtable}
